% SPEC-SH-004 --- Domain Normalizer

% Shared preamble for all spec documents.
% Include from each spec: % Shared preamble for all spec documents.
% Include from each spec: % Shared preamble for all spec documents.
% Include from each spec: \input{../preamble.tex} (adjust depth).

\documentclass[11pt,a4paper]{article}

\usepackage[utf8]{inputenc}
\usepackage[T1]{fontenc}
\usepackage{lmodern}
\usepackage[a4paper,margin=2.3cm]{geometry}
\usepackage{parskip}
\usepackage{microtype}
\usepackage[dvipsnames]{xcolor}
\usepackage{enumitem}
\usepackage{listings}
\usepackage{tcolorbox}
\usepackage{titlesec}
\usepackage{fancyhdr}
\usepackage{array}
\usepackage{longtable}
\usepackage{booktabs}
\usepackage{amsmath}
\usepackage{amssymb}
\usepackage{hyperref}
\usepackage{cleveref}

\tcbuselibrary{skins,breakable}

\definecolor{specBlue}{RGB}{30,64,132}
\definecolor{specGray}{RGB}{90,90,90}
\definecolor{specAccent}{RGB}{198,40,40}
\definecolor{specGreen}{RGB}{30,110,60}

\hypersetup{
  colorlinks=true,
  linkcolor=specBlue,
  urlcolor=specBlue,
  citecolor=specBlue,
  pdfborder={0 0 0}
}

\titleformat{\section}{\Large\bfseries\color{specBlue}}{\thesection}{0.75em}{}
\titleformat{\subsection}{\large\bfseries\color{specBlue!80}}{\thesubsection}{0.6em}{}

\makeatletter
\newcommand{\specID}[1]{\def\@specID{#1}}
\newcommand{\specTitle}[1]{\def\@specTitle{#1}}
\newcommand{\specStatus}[1]{\def\@specStatus{#1}}
\newcommand{\specVersion}[1]{\def\@specVersion{#1}}
\newcommand{\specOwner}[1]{\def\@specOwner{#1}}
\newcommand{\specTraces}[1]{\def\@specTraces{#1}}

\specID{SPEC-XXX}
\specTitle{Untitled Spec}
\specStatus{DRAFT}
\specVersion{0.1}
\specOwner{Jeremie Nunez}
\specTraces{}

\newcommand{\makespecheader}{%
  \begin{center}
    {\LARGE\bfseries\color{specBlue}\@specTitle}\\[0.4em]
    {\small\color{specGray}\texttt{\@specID} \quad v\@specVersion \quad \textsc{\@specStatus} \quad \@specOwner}
  \end{center}
  \vspace{0.4em}
  \hrule
  \vspace{0.8em}
}

\pagestyle{fancy}
\fancyhf{}
\fancyhead[L]{\small\color{specGray}\@specID}
\fancyhead[R]{\small\color{specGray}\@specTitle}
\fancyfoot[C]{\small\color{specGray}\thepage}
\renewcommand{\headrulewidth}{0pt}
\makeatother

\newtcolorbox{invariant}{
  colback=specBlue!5, colframe=specBlue!75,
  title=Invariant, fonttitle=\bfseries, breakable,
  boxrule=0.5pt, arc=2pt
}

\newtcolorbox{scenario}[1]{
  colback=white, colframe=specBlue!50,
  title={Scenario: #1}, fonttitle=\bfseries, breakable,
  boxrule=0.5pt, arc=2pt
}

\newtcolorbox{ac}[1]{
  colback=specAccent!5, colframe=specAccent!60,
  title={Acceptance Criterion #1}, fonttitle=\bfseries, breakable,
  boxrule=0.5pt, arc=2pt
}

\newtcolorbox{nongoal}{
  colback=specGray!8, colframe=specGray!60,
  title=Non-Goal, fonttitle=\bfseries, breakable,
  boxrule=0.5pt, arc=2pt
}

\lstdefinestyle{codestyle}{
  basicstyle=\ttfamily\small,
  breaklines=true,
  showstringspaces=false,
  frame=single,
  framesep=4pt,
  rulecolor=\color{specBlue!30},
  backgroundcolor=\color{specBlue!2},
  xleftmargin=0.5em,
  xrightmargin=0.5em,
  keywordstyle=\color{specBlue}\bfseries,
  commentstyle=\color{specGray}\itshape,
  stringstyle=\color{specGreen}
}
\lstset{
  inputencoding=utf8,
  extendedchars=true,
  literate=
    {—}{{--}}1
    {–}{{-}}1
    {…}{{...}}1
    {’}{{'}}1
    {‘}{{'}}1
    {“}{{``}}1
    {”}{{''}}1
    {é}{{\'e}}1
    {è}{{\`e}}1
    {ê}{{\^e}}1
    {à}{{\`a}}1
    {â}{{\^a}}1
    {ç}{{\c{c}}}1
    {ô}{{\^o}}1
    {→}{{$\to$}}1
    {←}{{$\leftarrow$}}1
    {×}{{$\times$}}1
    {≥}{{$\geq$}}1
    {≤}{{$\leq$}}1
    {≠}{{$\neq$}}1
    {∈}{{$\in$}}1
    {⌈}{{$\lceil$}}1
    {⌉}{{$\rceil$}}1
    {∞}{{$\infty$}}1
    {α}{{$\alpha$}}1
    {β}{{$\beta$}}1
}
\lstdefinelanguage{TypeScript}{
  keywords={const,let,var,function,return,if,else,for,while,class,interface,type,extends,implements,import,export,from,as,async,await,new,this,true,false,null,undefined,void,any,unknown,never,string,number,boolean,readonly,private,public,protected,enum},
  sensitive=true,
  comment=[l]{//},
  morecomment=[s]{/*}{*/},
  morestring=[b]",
  morestring=[b]',
  morestring=[b]`
}
\lstset{style=codestyle}

\newcommand{\Given}{\textbf{\color{specBlue}Given} }
\newcommand{\When}{\textbf{\color{specBlue}When} }
\newcommand{\Then}{\textbf{\color{specBlue}Then} }
\providecommand{\And}{}
\renewcommand{\And}{\textbf{\color{specBlue}And} }
 (adjust depth).

\documentclass[11pt,a4paper]{article}

\usepackage[utf8]{inputenc}
\usepackage[T1]{fontenc}
\usepackage{lmodern}
\usepackage[a4paper,margin=2.3cm]{geometry}
\usepackage{parskip}
\usepackage{microtype}
\usepackage[dvipsnames]{xcolor}
\usepackage{enumitem}
\usepackage{listings}
\usepackage{tcolorbox}
\usepackage{titlesec}
\usepackage{fancyhdr}
\usepackage{array}
\usepackage{longtable}
\usepackage{booktabs}
\usepackage{amsmath}
\usepackage{amssymb}
\usepackage{hyperref}
\usepackage{cleveref}

\tcbuselibrary{skins,breakable}

\definecolor{specBlue}{RGB}{30,64,132}
\definecolor{specGray}{RGB}{90,90,90}
\definecolor{specAccent}{RGB}{198,40,40}
\definecolor{specGreen}{RGB}{30,110,60}

\hypersetup{
  colorlinks=true,
  linkcolor=specBlue,
  urlcolor=specBlue,
  citecolor=specBlue,
  pdfborder={0 0 0}
}

\titleformat{\section}{\Large\bfseries\color{specBlue}}{\thesection}{0.75em}{}
\titleformat{\subsection}{\large\bfseries\color{specBlue!80}}{\thesubsection}{0.6em}{}

\makeatletter
\newcommand{\specID}[1]{\def\@specID{#1}}
\newcommand{\specTitle}[1]{\def\@specTitle{#1}}
\newcommand{\specStatus}[1]{\def\@specStatus{#1}}
\newcommand{\specVersion}[1]{\def\@specVersion{#1}}
\newcommand{\specOwner}[1]{\def\@specOwner{#1}}
\newcommand{\specTraces}[1]{\def\@specTraces{#1}}

\specID{SPEC-XXX}
\specTitle{Untitled Spec}
\specStatus{DRAFT}
\specVersion{0.1}
\specOwner{Jeremie Nunez}
\specTraces{}

\newcommand{\makespecheader}{%
  \begin{center}
    {\LARGE\bfseries\color{specBlue}\@specTitle}\\[0.4em]
    {\small\color{specGray}\texttt{\@specID} \quad v\@specVersion \quad \textsc{\@specStatus} \quad \@specOwner}
  \end{center}
  \vspace{0.4em}
  \hrule
  \vspace{0.8em}
}

\pagestyle{fancy}
\fancyhf{}
\fancyhead[L]{\small\color{specGray}\@specID}
\fancyhead[R]{\small\color{specGray}\@specTitle}
\fancyfoot[C]{\small\color{specGray}\thepage}
\renewcommand{\headrulewidth}{0pt}
\makeatother

\newtcolorbox{invariant}{
  colback=specBlue!5, colframe=specBlue!75,
  title=Invariant, fonttitle=\bfseries, breakable,
  boxrule=0.5pt, arc=2pt
}

\newtcolorbox{scenario}[1]{
  colback=white, colframe=specBlue!50,
  title={Scenario: #1}, fonttitle=\bfseries, breakable,
  boxrule=0.5pt, arc=2pt
}

\newtcolorbox{ac}[1]{
  colback=specAccent!5, colframe=specAccent!60,
  title={Acceptance Criterion #1}, fonttitle=\bfseries, breakable,
  boxrule=0.5pt, arc=2pt
}

\newtcolorbox{nongoal}{
  colback=specGray!8, colframe=specGray!60,
  title=Non-Goal, fonttitle=\bfseries, breakable,
  boxrule=0.5pt, arc=2pt
}

\lstdefinestyle{codestyle}{
  basicstyle=\ttfamily\small,
  breaklines=true,
  showstringspaces=false,
  frame=single,
  framesep=4pt,
  rulecolor=\color{specBlue!30},
  backgroundcolor=\color{specBlue!2},
  xleftmargin=0.5em,
  xrightmargin=0.5em,
  keywordstyle=\color{specBlue}\bfseries,
  commentstyle=\color{specGray}\itshape,
  stringstyle=\color{specGreen}
}
\lstset{
  inputencoding=utf8,
  extendedchars=true,
  literate=
    {—}{{--}}1
    {–}{{-}}1
    {…}{{...}}1
    {’}{{'}}1
    {‘}{{'}}1
    {“}{{``}}1
    {”}{{''}}1
    {é}{{\'e}}1
    {è}{{\`e}}1
    {ê}{{\^e}}1
    {à}{{\`a}}1
    {â}{{\^a}}1
    {ç}{{\c{c}}}1
    {ô}{{\^o}}1
    {→}{{$\to$}}1
    {←}{{$\leftarrow$}}1
    {×}{{$\times$}}1
    {≥}{{$\geq$}}1
    {≤}{{$\leq$}}1
    {≠}{{$\neq$}}1
    {∈}{{$\in$}}1
    {⌈}{{$\lceil$}}1
    {⌉}{{$\rceil$}}1
    {∞}{{$\infty$}}1
    {α}{{$\alpha$}}1
    {β}{{$\beta$}}1
}
\lstdefinelanguage{TypeScript}{
  keywords={const,let,var,function,return,if,else,for,while,class,interface,type,extends,implements,import,export,from,as,async,await,new,this,true,false,null,undefined,void,any,unknown,never,string,number,boolean,readonly,private,public,protected,enum},
  sensitive=true,
  comment=[l]{//},
  morecomment=[s]{/*}{*/},
  morestring=[b]",
  morestring=[b]',
  morestring=[b]`
}
\lstset{style=codestyle}

\newcommand{\Given}{\textbf{\color{specBlue}Given} }
\newcommand{\When}{\textbf{\color{specBlue}When} }
\newcommand{\Then}{\textbf{\color{specBlue}Then} }
\providecommand{\And}{}
\renewcommand{\And}{\textbf{\color{specBlue}And} }
 (adjust depth).

\documentclass[11pt,a4paper]{article}

\usepackage[utf8]{inputenc}
\usepackage[T1]{fontenc}
\usepackage{lmodern}
\usepackage[a4paper,margin=2.3cm]{geometry}
\usepackage{parskip}
\usepackage{microtype}
\usepackage[dvipsnames]{xcolor}
\usepackage{enumitem}
\usepackage{listings}
\usepackage{tcolorbox}
\usepackage{titlesec}
\usepackage{fancyhdr}
\usepackage{array}
\usepackage{longtable}
\usepackage{booktabs}
\usepackage{amsmath}
\usepackage{amssymb}
\usepackage{hyperref}
\usepackage{cleveref}

\tcbuselibrary{skins,breakable}

\definecolor{specBlue}{RGB}{30,64,132}
\definecolor{specGray}{RGB}{90,90,90}
\definecolor{specAccent}{RGB}{198,40,40}
\definecolor{specGreen}{RGB}{30,110,60}

\hypersetup{
  colorlinks=true,
  linkcolor=specBlue,
  urlcolor=specBlue,
  citecolor=specBlue,
  pdfborder={0 0 0}
}

\titleformat{\section}{\Large\bfseries\color{specBlue}}{\thesection}{0.75em}{}
\titleformat{\subsection}{\large\bfseries\color{specBlue!80}}{\thesubsection}{0.6em}{}

\makeatletter
\newcommand{\specID}[1]{\def\@specID{#1}}
\newcommand{\specTitle}[1]{\def\@specTitle{#1}}
\newcommand{\specStatus}[1]{\def\@specStatus{#1}}
\newcommand{\specVersion}[1]{\def\@specVersion{#1}}
\newcommand{\specOwner}[1]{\def\@specOwner{#1}}
\newcommand{\specTraces}[1]{\def\@specTraces{#1}}

\specID{SPEC-XXX}
\specTitle{Untitled Spec}
\specStatus{DRAFT}
\specVersion{0.1}
\specOwner{Jeremie Nunez}
\specTraces{}

\newcommand{\makespecheader}{%
  \begin{center}
    {\LARGE\bfseries\color{specBlue}\@specTitle}\\[0.4em]
    {\small\color{specGray}\texttt{\@specID} \quad v\@specVersion \quad \textsc{\@specStatus} \quad \@specOwner}
  \end{center}
  \vspace{0.4em}
  \hrule
  \vspace{0.8em}
}

\pagestyle{fancy}
\fancyhf{}
\fancyhead[L]{\small\color{specGray}\@specID}
\fancyhead[R]{\small\color{specGray}\@specTitle}
\fancyfoot[C]{\small\color{specGray}\thepage}
\renewcommand{\headrulewidth}{0pt}
\makeatother

\newtcolorbox{invariant}{
  colback=specBlue!5, colframe=specBlue!75,
  title=Invariant, fonttitle=\bfseries, breakable,
  boxrule=0.5pt, arc=2pt
}

\newtcolorbox{scenario}[1]{
  colback=white, colframe=specBlue!50,
  title={Scenario: #1}, fonttitle=\bfseries, breakable,
  boxrule=0.5pt, arc=2pt
}

\newtcolorbox{ac}[1]{
  colback=specAccent!5, colframe=specAccent!60,
  title={Acceptance Criterion #1}, fonttitle=\bfseries, breakable,
  boxrule=0.5pt, arc=2pt
}

\newtcolorbox{nongoal}{
  colback=specGray!8, colframe=specGray!60,
  title=Non-Goal, fonttitle=\bfseries, breakable,
  boxrule=0.5pt, arc=2pt
}

\lstdefinestyle{codestyle}{
  basicstyle=\ttfamily\small,
  breaklines=true,
  showstringspaces=false,
  frame=single,
  framesep=4pt,
  rulecolor=\color{specBlue!30},
  backgroundcolor=\color{specBlue!2},
  xleftmargin=0.5em,
  xrightmargin=0.5em,
  keywordstyle=\color{specBlue}\bfseries,
  commentstyle=\color{specGray}\itshape,
  stringstyle=\color{specGreen}
}
\lstset{
  inputencoding=utf8,
  extendedchars=true,
  literate=
    {—}{{--}}1
    {–}{{-}}1
    {…}{{...}}1
    {’}{{'}}1
    {‘}{{'}}1
    {“}{{``}}1
    {”}{{''}}1
    {é}{{\'e}}1
    {è}{{\`e}}1
    {ê}{{\^e}}1
    {à}{{\`a}}1
    {â}{{\^a}}1
    {ç}{{\c{c}}}1
    {ô}{{\^o}}1
    {→}{{$\to$}}1
    {←}{{$\leftarrow$}}1
    {×}{{$\times$}}1
    {≥}{{$\geq$}}1
    {≤}{{$\leq$}}1
    {≠}{{$\neq$}}1
    {∈}{{$\in$}}1
    {⌈}{{$\lceil$}}1
    {⌉}{{$\rceil$}}1
    {∞}{{$\infty$}}1
    {α}{{$\alpha$}}1
    {β}{{$\beta$}}1
}
\lstdefinelanguage{TypeScript}{
  keywords={const,let,var,function,return,if,else,for,while,class,interface,type,extends,implements,import,export,from,as,async,await,new,this,true,false,null,undefined,void,any,unknown,never,string,number,boolean,readonly,private,public,protected,enum},
  sensitive=true,
  comment=[l]{//},
  morecomment=[s]{/*}{*/},
  morestring=[b]",
  morestring=[b]',
  morestring=[b]`
}
\lstset{style=codestyle}

\newcommand{\Given}{\textbf{\color{specBlue}Given} }
\newcommand{\When}{\textbf{\color{specBlue}When} }
\newcommand{\Then}{\textbf{\color{specBlue}Then} }
\providecommand{\And}{}
\renewcommand{\And}{\textbf{\color{specBlue}And} }


\specID{SPEC-SH-004}
\specTitle{Domain Normalizer --- canonical form derivation for noisy identifiers}
\specStatus{DRAFT}
\specVersion{0.1}
\specOwner{Jeremie Nunez}

\begin{document}
\makespecheader

\section{Context}
Records arrive from heterogeneous sources using inconsistent spellings for the same domain identifier: mixed case, surrounding whitespace, separator variants (space vs.\ hyphen vs.\ underscore), diacritics, stray punctuation, and occasional prefixes or suffixes that are semantically noise. Downstream indexing, deduplication, joins and embeddings all require a single \emph{canonical form} per logical identifier. The domain normalizer is the shared utility that derives this canonical form deterministically from any raw input, so that equality of canonical forms is a reliable proxy for logical equality.

\section{Goals}
\begin{itemize}[leftmargin=1.2em]
  \item Produce a deterministic, side-effect-free canonical string from any input string.
  \item Collapse the common noise axes (case, whitespace, separators, diacritics) without destroying semantic content.
  \item Preserve reversibility metadata: callers can inspect what was stripped or folded, and recover the original input.
  \item Be safe on pathological inputs (empty string, only punctuation, very long strings, mixed scripts) without throwing.
\end{itemize}

\section{Non-Goals}
\begin{nongoal}
The normalizer does \emph{not} perform fuzzy matching, phonetic encoding (Soundex / Metaphone), translation between languages, or synonym resolution. It does not know any domain ontology: it operates on the string surface. Mapping canonical forms to business entities is the job of a higher-level resolver.
\end{nongoal}

\section{Invariants}
\begin{invariant}
Idempotence: for any input \(s\), \(\mathrm{normalize}(\mathrm{normalize}(s)) = \mathrm{normalize}(s)\).
\end{invariant}

\begin{invariant}
Determinism: \(\mathrm{normalize}(s)\) depends only on \(s\) and on compile-time configuration; no locale, time, or process state influences the output.
\end{invariant}

\begin{invariant}
Non-destruction: if two inputs differ in a character that is neither whitespace, nor a recognised separator, nor a diacritic, nor case, then their canonical forms differ.
\end{invariant}

\section{Interface}

\begin{lstlisting}[language=TypeScript]
export interface NormalizeOptions {
  readonly separator?: "-" | "_" | "";   // default: "-"
  readonly stripDiacritics?: boolean;    // default: true
  readonly lowercase?: boolean;          // default: true
  readonly collapseWhitespace?: boolean; // default: true
}

export interface NormalizeResult {
  readonly canonical: string;
  readonly original: string;
  readonly strippedDiacritics: boolean;
  readonly collapsedSeparators: boolean;
}

export function normalizeDomain(
  input: string,
  options?: NormalizeOptions
): NormalizeResult;

export function canonical(
  input: string,
  options?: NormalizeOptions
): string; // shorthand: returns result.canonical
\end{lstlisting}

\section{Scenarios}

\begin{scenario}{Case, whitespace, diacritics all folded}
\Given the input \texttt{"  C\^{o}te-Rotie  "} \\
\When \texttt{normalizeDomain} is called with defaults \\
\Then \texttt{canonical} equals \texttt{"cote-rotie"} \\
\And \texttt{strippedDiacritics} is \texttt{true}.
\end{scenario}

\begin{scenario}{Separator variants collapse to one}
\Given two inputs \texttt{"foo bar"}, \texttt{"foo\_bar"}, \texttt{"foo--bar"} \\
\When each is normalised with defaults \\
\Then all three produce \texttt{canonical} equal to \texttt{"foo-bar"}.
\end{scenario}

\begin{scenario}{Idempotence}
\Given any input \(s\) \\
\When normalization is applied twice \\
\Then the second application returns the same \texttt{canonical} string as the first, and the \texttt{original} field still reflects the first-stage input.
\end{scenario}

\begin{scenario}{Empty or punctuation-only input}
\Given inputs \texttt{""}, \texttt{"   "}, \texttt{"---"} \\
\When normalization runs \\
\Then \texttt{canonical} equals \texttt{""} \\
\And no exception is thrown.
\end{scenario}

\begin{scenario}{Non-Latin script preserved}
\Given input \texttt{"\^{E}l\'{e}gant 2024"} with diacritic stripping on \\
\When normalised \\
\Then \texttt{canonical} equals \texttt{"elegant-2024"}, preserving the digit run as a single token.
\end{scenario}

\section{Acceptance Criteria}

\begin{ac}{AC-1}
Default normalization lowercases, strips leading / trailing whitespace, folds diacritics on Latin characters, and replaces any run of whitespace, hyphen, or underscore by a single configured separator (default \texttt{"-"}).
\end{ac}

\begin{ac}{AC-2}
For every string \(s\), \texttt{canonical(canonical(s)) === canonical(s)} (idempotence).
\end{ac}

\begin{ac}{AC-3}
Two inputs that differ only in case, whitespace, or separator variant produce identical canonical forms under default options.
\end{ac}

\begin{ac}{AC-4}
Two inputs that differ in a non-folded character produce different canonical forms.
\end{ac}

\begin{ac}{AC-5}
\texttt{normalizeDomain("")} and \texttt{normalizeDomain("   ")} both return \texttt{canonical === ""} without throwing.
\end{ac}

\begin{ac}{AC-6}
When \texttt{stripDiacritics} is \texttt{false}, diacritics are preserved verbatim in the canonical form.
\end{ac}

\begin{ac}{AC-7}
The \texttt{original} field on the result always equals the input passed in, regardless of options.
\end{ac}

\begin{ac}{AC-8}
On inputs up to \(10^{5}\) characters, normalization completes in under \(10\) ms on the reference CI runner and allocates \(O(n)\) memory.
\end{ac}

\section{Traceability}

\begin{center}
\begin{tabular}{@{}lll@{}}
\toprule
AC & Test file & Test name \\
\midrule
AC-1 & \texttt{packages/shared/tests/domain-normalizer.spec.ts} & \texttt{default options fold case, whitespace, diacritics} \\
AC-2 & \texttt{packages/shared/tests/domain-normalizer.spec.ts} & \texttt{normalization is idempotent} \\
AC-3 & \texttt{packages/shared/tests/domain-normalizer.spec.ts} & \texttt{separator variants collapse to canonical separator} \\
AC-4 & \texttt{packages/shared/tests/domain-normalizer.spec.ts} & \texttt{non-folded characters remain distinguishing} \\
AC-5 & \texttt{packages/shared/tests/domain-normalizer.spec.ts} & \texttt{empty and whitespace-only inputs return empty canonical} \\
AC-6 & \texttt{packages/shared/tests/domain-normalizer.spec.ts} & \texttt{stripDiacritics=false preserves accents} \\
AC-7 & \texttt{packages/shared/tests/domain-normalizer.spec.ts} & \texttt{original field round-trips the input} \\
AC-8 & \texttt{packages/shared/tests/domain-normalizer.bench.ts} & \texttt{normalizes 100k-character input under 10ms} \\
\bottomrule
\end{tabular}
\end{center}

\section{Open Questions}
\begin{itemize}[leftmargin=1.2em]
  \item Should non-Latin diacritic-bearing scripts (e.g.\ Cyrillic, Greek) be folded by the same NFD+strip pass, or preserved?
  \item Is an optional \texttt{maxLength} truncation parameter in scope, or left to callers?
  \item Should the API expose a streaming variant for very large bulk ingestion, or is the synchronous signature sufficient?
\end{itemize}

\end{document}
